\documentclass[10pt]{scrartcl}
\usepackage[mdthm,colorsec,noauthor,noasy]{evan}
\usepackage{handout}
\usepackage{figure-ink}
\newcommand{\aug}[2]{\left[\begin{array}{@{}#1@{}}#2\end{array}\right]}
\newcommand{\vv}[1]{\mathbf{#1}}
\DeclareMathOperator{\spn}{Span}\DeclareMathOperator{\Nul}{Nul}\DeclareMathOperator{\Col}{Col}\DeclareMathOperator{\Row}{Row}\DeclareMathOperator{\tr}{tr}
\title{Math 54 --- Week 7}
\subtitle{Coordinate systems; change of basis; eigenvectors and eigenvalues; the characteristic equation (Lay 4.4, 4.6, 5.1--5.2)}
\author{}\date{}\ihead{}
\begin{document}
\maketitle

\section{Facts used from Weeks 1--6}

\begin{enumerate}[label=(W\arabic*),leftmargin=2.8em,itemsep=4pt]
  \ii \target{fact-w1}\textbf{Bases (\xlink{math54}{2026-09-28-bases-dimension-rank}{thm-lay-4-12}{Week~6}):} A basis of $V$ is a linearly independent spanning set. Every basis of $V$ has $n=\dim V$ vectors; in an $n$-dimensional space, any $n$ independent vectors form a basis, and more than $n$ vectors are dependent (\xlink{math54}{2026-09-28-bases-dimension-rank}{thm-lay-4-9}{Lay 4.5}).

  \ii \target{fact-w2}\textbf{Invertible matrix theorem (\xlink{math54}{2026-09-14-matrix-algebra-inverses}{thm-lay-2-8}{Weeks~4--6}):} For square $A$, $A$ is invertible $\iff A\vv x=\vv 0$ has only the trivial solution $\iff \det A\ne 0 \iff$ the columns of $A$ form a basis of $\RR^n$. For $2\times 2$ matrices,
  \[
    \begin{bmatrix}a&b\\c&d\end{bmatrix}^{-1}=\frac{1}{ad-bc}\begin{bmatrix}d&-b\\-c&a\end{bmatrix}.
  \]

  \ii \target{fact-w3}\textbf{Determinants (\xlink{math54}{2026-09-21-determinants-vector-spaces}{thm-lay-3-6}{Weeks~4--5}):} $\det(AB)=\det A\det B$, $\det(A^{\mathsf T})=\det A$, row operations modify $\det$ predictably, and the determinant of a triangular matrix is the product of its diagonal entries.

  \ii \target{fact-w4}\textbf{Subspaces and row equivalence (\xlink{math54}{2026-09-28-bases-dimension-rank}{thm-lay-4-6}{Weeks~5--6}):} $\Nul A$ is a subspace of $\RR^n$, with a basis given by the parametric vector form of $A\vv x=\vv 0$. Row operations preserve linear relations among the columns of a matrix.

  \ii \target{fact-w5}\textbf{One-to-one linear transformations (\xlink{math54}{2026-09-07-independence-linear-transformations}{thm-lay-1-11}{Lay 1.9 Thm.~11}):} A linear map $T\colon V\to W$ between any vector spaces is one-to-one if and only if $T(\vv x)=\vv 0$ holds only for $\vv x=\vv 0$. The Week~3 proof uses only linearity: $T(\vv 0)=\vv 0$, and $T(\vv u)=T(\vv v)$ gives $T(\vv u-\vv v)=\vv 0$, so $\vv u-\vv v=\vv 0$.

  \ii \target{fact-w6}\textbf{Linear dependence lemma (\xlink{math54}{2026-09-07-independence-linear-transformations}{thm-lay-1-7}{Lay 1.7 Thm.~7}, \xlink{math54}{2026-09-28-bases-dimension-rank}{thm-lay-4-4}{Lay 4.3 Thm.~4}):} If $\{\vv v_1,\dots,\vv v_p\}$ is dependent and $\vv v_1\ne\vv 0$, then some $\vv v_j$ ($j>1$) is a linear combination of $\vv v_1,\dots,\vv v_{j-1}$.
\end{enumerate}

\section{This week}

\subsection{Coordinate systems (Lay 4.4)}

\begin{qbox}
  Given the slanted basis $\mathcal B=\{\vv b_1,\vv b_2\}=\{(2,1),(-1,1)\}$ of $\RR^2$, how do we describe $\vv x=(4,5)$ by a coordinate list relative to $\mathcal B$, why is this coordinate list \emph{unique}, and how does a coordinate system turn polynomials or matrices into ordinary column vectors in $\RR^n$?
\end{qbox}

In standard Euclidean coordinates we measure position using grid steps along $(1,0)$ and $(0,1)$. Because \emph{any} basis $\mathcal B$ spans $V$ with zero redundancy, $\mathcal B$ sets up a new skewed coordinate grid on $V$ in which every vector has a unique address.

\begin{theorem*}[The Unique Representation Theorem (Lay 4.4, Thm.~8)]\target{thm-lay-4-8}
  Let $\mathcal B=\{\vv b_1,\dots,\vv b_n\}$ be a basis for a vector space $V$. Then for each $\vv x\in V$, there exists a \textbf{unique} set of scalars $c_1,\dots,c_n$ such that
  \[
    \vv x = c_1\vv b_1 + \dots + c_n\vv b_n.
  \]
\end{theorem*}
\begin{ideabox}
  Spanning gives \emph{existence} of weights $c_1,\dots,c_n$, while linear independence gives \emph{uniqueness}: subtracting two representations produces a linear combination equal to $\vv 0$, forcing all differences $c_i-d_i$ to vanish.
\end{ideabox}
\begin{proof}
  \begin{steps}
    \ii Since $\mathcal B$ spans $V$, at least one representation $\vv x=c_1\vv b_1+\dots+c_n\vv b_n$ exists.
    \ii If also $\vv x=d_1\vv b_1+\dots+d_n\vv b_n$, subtracting yields $\vv 0=(c_1-d_1)\vv b_1+\dots+(c_n-d_n)\vv b_n$. By linear independence of $\mathcal B$, $c_i-d_i=0$ for every $i$, so $c_i=d_i$.
  \end{steps}
\end{proof}

\noindent\textbf{Definition.} The unique weights $c_1,\dots,c_n$ are the \vocab{coordinates of $\vv x$ relative to $\mathcal B$}, and the column vector $[\vv x]_{\mathcal B}=(c_1,\dots,c_n)\in\RR^n$ is the \vocab{$\mathcal B$-coordinate vector of $\vv x$}. The assignment $\vv x\mapsto[\vv x]_{\mathcal B}$ is the \vocab{coordinate mapping}.

\medskip\noindent\textbf{Coordinates in $\RR^n$ and the matrix $P_{\mathcal B}$.} Putting the basis vectors as columns forms $P_{\mathcal B}=[\,\vv b_1\ \cdots\ \vv b_n\,]$, the \vocab{change-of-coordinates matrix} from $\mathcal B$ to the standard basis $\mathcal E$. By definition of matrix--vector multiplication,
\[
  \vv x = c_1\vv b_1+\dots+c_n\vv b_n = P_{\mathcal B}[\vv x]_{\mathcal B},
  \qquad\text{hence}\qquad
  [\vv x]_{\mathcal B} = P_{\mathcal B}^{-1}\vv x.
\]
Here $P_{\mathcal B}$ is guaranteed invertible because its columns form a basis of $\RR^n$ (\ilink{fact-w2}{W2}). In practice, rather than inverting $P_{\mathcal B}$, row reduce the augmented matrix $[\,\vv b_1\ \cdots\ \vv b_n\mid \vv x\,]\sim [\,I\mid [\vv x]_{\mathcal B}\,]$.

\revisit Solving $c_1(2,1)+c_2(-1,1)=(4,5)$ (swapping rows first to avoid fractions):
\[
  \aug{cc|c}{2&-1&4\\1&1&5}
  \xrightarrow{R_1\leftrightarrow R_2}
  \aug{cc|c}{1&1&5\\2&-1&4}
  \xrightarrow{R_2-2R_1}
  \aug{cc|c}{1&1&5\\0&-3&-6}
  \xrightarrow{-\frac13 R_2,\ R_1-R_2}
  \aug{cc|c}{1&0&3\\0&1&2},
  \qquad [\vv x]_{\mathcal B} = \begin{bmatrix}3\\2\end{bmatrix}.
\]
Check: $3(2,1)+2(-1,1)=(4,5)$. These are the only weights, by \ilink{thm-lay-4-8}{Thm.~4.8}.

\begin{fig}
\begin{tikzpicture}[>=Stealth]
  \begin{scope}[scale=0.56]
    % Soft borderless gray wash for the fundamental basis cell [b1, b2]
    \path[pattern=inkgrain,pattern color=inkblue] (0,0) -- (2,1) -- (1,2) -- (-1,1) -- cycle;

    % Delicate Spivak skewed B-grid for c1 in [-1,3], c2 in [0,2]
    \foreach \i in {-1,0,1,2,3} {
      \InkStroke[black!45,line width=0.4pt]{({2*\i},{\i}) -- ({2*\i-2},{\i+2})}
    }
    \foreach \j in {0,1,2} {
      \InkStroke[black!45,line width=0.4pt]{({-2-\j},{-1+\j}) -- ({6-\j},{3+\j})}
    }

    % Hairline axes
    \draw[ink line,figaxis] (-4.4,0) -- (7.3,0) node[right,font=\scriptsize] {\InkLabel{$x_1$}};
    \draw[ink line,figaxis] (0,-1.4) -- (0,6.1) node[above,font=\scriptsize] {\InkLabel{$x_2$}};

    % Basis vectors b1, b2 (restrained dark ink)
    \InkStroke[figvec,draw=black,line width=0.65pt]{(0,0) -- (2,1)}
    \node[font=\scriptsize] at (2.15,0.42) {\InkLabel{$\vv b_1$}};

    \InkStroke[figvec,draw=black,line width=0.65pt]{(0,0) -- (-1,1)}
    \node[font=\scriptsize] at (-1.0,1.8) {\InkLabel{$\vv b_2$}};

    % Delicate dotted coordinate path 3 b1 + 2 b2
    \draw[ink guide arrow,draw=black] (2,1) -- (6,3) node[below right=2pt,font=\scriptsize,text=black] {\InkLabel{$3\vv b_1$}};
    \draw[ink guide arrow,draw=black,shorten >=3pt] (6,3) -- (4,5) node[midway,above right=2pt,font=\scriptsize,text=black] {\InkLabel{$2\vv b_2$}};

    % Focal vector x
    \InkStroke[figvec,draw=black,line width=0.65pt,shorten >=2.8pt]{(0,0) -- (4,5)}
    \node[figdot] at (4,5) {};
    \node[above=4pt,font=\scriptsize] at (4,5) {\InkLabel{$\vv x = 3\vv b_1+2\vv b_2$}};
  \end{scope}
\end{tikzpicture}
\figcap{The grid of $\mathcal B=\{(2,1),(-1,1)\}$: $3$ steps along $\vv b_1$ and $2$ steps along $\vv b_2$ reach $\vv x=(4,5)$, so $[\vv x]_{\mathcal B}=(3,2)$.}
\end{fig}

\begin{theorem*}[The Coordinate Mapping Isomorphism (Lay 4.4, Thm.~9)]\target{thm-lay-4-9}
  Let $\mathcal B=\{\vv b_1,\dots,\vv b_n\}$ be a basis for a vector space $V$. Then the coordinate mapping $\vv x\mapsto[\vv x]_{\mathcal B}$ is a one-to-one linear transformation from $V$ onto $\RR^n$.
\end{theorem*}
\begin{proof}
  \begin{steps}
    \ii \emph{Linearity.} Write $\vv u=\sum_{i=1}^n c_i\vv b_i$ and $\vv w=\sum_{i=1}^n d_i\vv b_i$. Then $\vv u+\vv w=\sum_{i=1}^n(c_i+d_i)\vv b_i$ and $r\vv u=\sum_{i=1}^n(rc_i)\vv b_i$. By uniqueness (\ilink{thm-lay-4-8}{Thm.~4.8}), $[\vv u+\vv w]_{\mathcal B}=[\vv u]_{\mathcal B}+[\vv w]_{\mathcal B}$ and $[r\vv u]_{\mathcal B}=r[\vv u]_{\mathcal B}$.
    \ii \emph{One-to-one.} If $[\vv x]_{\mathcal B}=\vv 0$, then $\vv x=0\vv b_1+\dots+0\vv b_n=\vv 0$, so by (\ilink{fact-w5}{W5}) the map is one-to-one.
    \ii \emph{Onto.} Any $(c_1,\dots,c_n)\in\RR^n$ is the coordinate vector of $\vv x=c_1\vv b_1+\dots+c_n\vv b_n\in V$.
  \end{steps}
\end{proof}

A one-to-one linear map from $V$ onto $\RR^n$ is an \vocab{isomorphism}: every $n$-dimensional real vector space $V$ behaves identically to $\RR^n$ once a basis is chosen. Because $\vv x\mapsto[\vv x]_{\mathcal B}$ is linear and $[\vv x]_{\mathcal B}=\vv 0\iff \vv x=\vv 0$, coordinates preserve \emph{all} linear relations:
\[
  c_1\vv u_1+\dots+c_p\vv u_p = \vv 0 \iff c_1[\vv u_1]_{\mathcal B}+\dots+c_p[\vv u_p]_{\mathcal B}=\vv 0.
\]
Thus questions about span and linear independence in $\mathbb P_n$ or $M_{2\times2}$ translate directly into columns in $\RR^N$.

\method{How to find coordinates, and how to recover a vector from its coordinates.}
\begin{steps}
  \ii \emph{From coordinates to the vector:} $\vv x=c_1\vv b_1+\dots+c_n\vv b_n$; in $\RR^n$ this is $\vv x=P_{\mathcal B}[\vv x]_{\mathcal B}$.
  \ii \emph{From the vector to coordinates in $\RR^n$:} row reduce $[\,\vv b_1\ \cdots\ \vv b_n\mid\vv x\,]$ to $[\,I\mid[\vv x]_{\mathcal B}\,]$ (or compute $P_{\mathcal B}^{-1}\vv x$ for a $2\times2$ matrix).
  \ii \emph{In $\mathbb P_n$ or another space:} write $c_1\vv b_1+\dots+c_n\vv b_n=\vv x$, expand, and equate the coefficients of $1,t,t^2,\dots$ (or the matrix entries); solve the resulting system.
  \ii \emph{Span and independence questions in $\mathbb P_n$:} replace each polynomial by its coordinate vector in $\{1,t,\dots,t^n\}$ and answer the question for those columns (\ilink{thm-lay-4-9}{Thm.~4.9}).
\end{steps}
\emph{Edge cases.} For a basis of a subspace $H$ of $\RR^n$ with $k$ vectors, $[\vv x]_{\mathcal B}$ has $k$ entries, and $\vv x\in H$ only if the system is consistent. The order of the basis vectors fixes the order of the coordinates.

\begin{example*}[A vector from its coordinates]
For $\mathcal B=\{(1,1),(-1,2)\}$ and $[\vv x]_{\mathcal B}=(3,1)$, $\vv x=3(1,1)+1(-1,2)=(2,5)$.
\end{example*}

\begin{example*}[Dependence of polynomials through coordinates]
Are $p_1=1+2t^2$, $p_2=4+t+5t^2$, $p_3=3+2t$ independent? Taking coordinate vectors in the standard basis $\{1,t,t^2\}$ gives columns $(1,0,2)^{\mathsf T}$, $(4,1,5)^{\mathsf T}$, $(3,2,0)^{\mathsf T}$:
    \[
      \begin{bmatrix}1&4&3\\0&1&2\\2&5&0\end{bmatrix}
      \xrightarrow{R_3-2R_1}
      \begin{bmatrix}1&4&3\\0&1&2\\0&-3&-6\end{bmatrix}
      \xrightarrow{R_3+3R_2,\ R_1-4R_2}
      \begin{bmatrix}1&0&-5\\0&1&2\\0&0&0\end{bmatrix}.
    \]
Since column~$3 = -5(\text{col }1)+2(\text{col }2)$ and row operations keep column relations (\ilink{fact-w4}{W4}), $p_3=-5p_1+2p_2$: \emph{dependent}.
\end{example*}

\begin{example*}[Coordinates relative to a basis of a plane in $\RR^3$]
Let $\vv v_1=(1,2,0)$, $\vv v_2=(0,1,1)$, $H=\spn\{\vv v_1,\vv v_2\}\subseteq\RR^3$ with basis $\mathcal B=\{\vv v_1,\vv v_2\}$. To check if $\vv x=(2,7,3)\in H$ and find $[\vv x]_{\mathcal B}$:
    \[
      \aug{cc|c}{1&0&2\\2&1&7\\0&1&3}
      \xrightarrow{R_2-2R_1,\ R_3-R_2}
      \aug{cc|c}{1&0&2\\0&1&3\\0&0&0}.
    \]
The system is consistent; thus $\vv x\in H$ with $[\vv x]_{\mathcal B}=(2,3)\in\RR^2$ (note $\dim H=2$, so $[\vv x]_{\mathcal B}\in\RR^2$ even though $\vv x\in\RR^3$).
\end{example*}

\begin{shortcuts}
  \ii \textbf{Coordinates in $\RR^n$:} Row reduce $[\,\vv b_1\ \cdots\ \vv b_n\mid \vv x\,]$ to find $[\vv x]_{\mathcal B}$.
  \ii \textbf{Span/independence in $\mathbb P_n$ or matrix spaces:} Convert to standard coordinate vectors in $\RR^N$ and row reduce (\ilink{thm-lay-4-9}{Thm.~4.9}).
  \ii \textbf{Polynomial in a non-standard basis:} Equate coefficients of $1,t,t^2,\dots$ (e.g.\ in $\mathbb P_1$ with $\mathcal B=\{1+t,1-t\}$, $c_1(1+t)+c_2(1-t)=5+t$ gives $c_1+c_2=5$ and $c_1-c_2=1$, so $[5+t]_{\mathcal B}=(3,2)$).
\end{shortcuts}

\begin{trapbox}
  \ii \textbf{Order matters:} Swapping $\mathcal B'=\{\vv b_2,\vv b_1\}$ swaps the coordinates from $(3,2)$ to $(2,3)$.
  \ii \textbf{Polynomial basis order:} With $\{1,t,t^2\}$, $1+2t^2$ has coordinates $(1,0,2)$, not $(2,0,1)$.
  \ii \textbf{$P_{\mathcal B}$ goes $\mathcal B\to\mathcal E$, not $\mathcal E\to\mathcal B$:} For the question, $P_{\mathcal B}(4,5)=(3,9)\neq[\vv x]_{\mathcal B}=(3,2)$; always use $[\vv x]_{\mathcal B}=P_{\mathcal B}^{-1}\vv x$.
  \ii \textbf{Subspace coordinates live in $\RR^{\dim H}$:} Not in the ambient $\RR^n$.
\end{trapbox}

\subsection{Change of basis (Lay 4.6)}

\begin{qbox}
  Let $\mathcal B=\{(1,1),(1,2)\}$ and $\mathcal C=\{(1,0),(2,1)\}$ be two bases of $\RR^2$. Suppose a vector $\vv x$ has $\mathcal B$-coordinates $[\vv x]_{\mathcal B}=(2,-1)$. How do we convert directly from $[\vv x]_{\mathcal B}$ to $[\vv x]_{\mathcal C}$ with a single matrix multiply?
\end{qbox}

\begin{theorem*}[Change of Coordinates (Lay 4.6, Thm.~15)]\target{thm-lay-4-15}
  Let $\mathcal B=\{\vv b_1,\dots,\vv b_n\}$ and $\mathcal C=\{\vv c_1,\dots,\vv c_n\}$ be bases of $V$. There exists a \textbf{unique} $n\times n$ matrix $P_{\mathcal C\leftarrow\mathcal B}$ such that
  \[
    [\vv x]_{\mathcal C} = P_{\mathcal C\leftarrow\mathcal B}[\vv x]_{\mathcal B} \qquad \text{for every }\vv x\in V.
  \]
  Its columns are the $\mathcal C$-coordinate vectors of the old basis $\mathcal B$:
  \[
    P_{\mathcal C\leftarrow\mathcal B} = \bigl[\,[\vv b_1]_{\mathcal C}\ \ [\vv b_2]_{\mathcal C}\ \ \cdots\ \ [\vv b_n]_{\mathcal C}\,\bigr].
  \]
  Moreover $P_{\mathcal C\leftarrow\mathcal B}$ is invertible, and $(P_{\mathcal C\leftarrow\mathcal B})^{-1} = P_{\mathcal B\leftarrow\mathcal C}$.
\end{theorem*}
\begin{ideabox}
  Since $\vv x\mapsto[\vv x]_{\mathcal C}$ is linear (\ilink{thm-lay-4-9}{Thm.~4.9}), knowing $[\vv x]_{\mathcal C}$ for any $\vv x=\sum c_j\vv b_j$ requires only knowing $[\vv b_j]_{\mathcal C}$ for each old basis vector $\vv b_j$: linearity factors out the scalars $c_j=([\vv x]_{\mathcal B})_j$.
\end{ideabox}
\begin{proof}
  \begin{steps}
    \ii Write $\vv x=\sum_{j=1}^n c_j\vv b_j$, so $[\vv x]_{\mathcal B}=(c_1,\dots,c_n)^{\mathsf T}$. By linearity of $[\cdot]_{\mathcal C}$ (\ilink{thm-lay-4-9}{Thm.~4.9}),
    \[
      [\vv x]_{\mathcal C} = \sum_{j=1}^n c_j[\vv b_j]_{\mathcal C} = P_{\mathcal C\leftarrow\mathcal B}[\vv x]_{\mathcal B}.
    \]
    \ii \emph{Uniqueness:} Applying any valid matrix $M$ to $\vv x=\vv b_j$ ($[\vv b_j]_{\mathcal B}=\vv e_j$) forces column $j$ of $M$ to equal $M\vv e_j = [\vv b_j]_{\mathcal C}$.
    \ii \emph{Invertibility \& inverse:} If $\sum d_j[\vv b_j]_{\mathcal C}=\vv 0$, then $[\sum d_j\vv b_j]_{\mathcal C}=\vv 0$, forcing $\sum d_j\vv b_j=\vv 0$ (\ilink{thm-lay-4-9}{Thm.~4.9}) and every $d_j=0$; hence columns are independent and $P_{\mathcal C\leftarrow\mathcal B}$ is invertible (\ilink{fact-w2}{W2}). Applying $(P_{\mathcal C\leftarrow\mathcal B})^{-1}[\vv x]_{\mathcal C}=[\vv x]_{\mathcal B}$ to $\vv x=\vv c_j$ shows column $j$ of $(P_{\mathcal C\leftarrow\mathcal B})^{-1}$ is $[\vv c_j]_{\mathcal B}$, i.e.\ $P_{\mathcal B\leftarrow\mathcal C}$.
  \end{steps}
\end{proof}

\method{How to find a change-of-coordinates matrix.}
\begin{steps}
  \ii Column $j$ of $P_{\mathcal C\leftarrow\mathcal B}$ is $[\vv b_j]_{\mathcal C}$: write each \emph{old} basis vector in terms of the \emph{new} basis.
  \ii In $\RR^n$, do all columns at once: row reduce $[\,\vv c_1\ \cdots\ \vv c_n\mid\vv b_1\ \cdots\ \vv b_n\,]$ to $[\,I\mid P_{\mathcal C\leftarrow\mathcal B}\,]$ (derivation below).
  \ii If $\mathcal C$ is the standard basis, $P_{\mathcal C\leftarrow\mathcal B}=P_{\mathcal B}$; no computation is needed.
  \ii For the other direction, $P_{\mathcal B\leftarrow\mathcal C}=(P_{\mathcal C\leftarrow\mathcal B})^{-1}$ (\ilink{thm-lay-4-15}{Thm.~4.15}).
  \ii Use it: $[\vv x]_{\mathcal C}=P_{\mathcal C\leftarrow\mathcal B}[\vv x]_{\mathcal B}$.
\end{steps}
\emph{Edge case.} If you are told $\vv b_1=2\vv c_1-\vv c_2$, the coefficients $(2,-1)$ form \emph{column}~1, not row~1.

\noindent\textbf{Computing $P_{\mathcal C\leftarrow\mathcal B}$ in $\RR^n$.} Column $j$ solves $[\vv c_1\ \cdots\ \vv c_n]\vv y=\vv b_j$. Solving all $n$ columns simultaneously reduces one augmented matrix:
\[
  [\,\vv c_1\ \cdots\ \vv c_n \mid \vv b_1\ \cdots\ \vv b_n\,] \longrightarrow [\,I \mid P_{\mathcal C\leftarrow\mathcal B}\,],
  \qquad\text{equivalently}\qquad P_{\mathcal C\leftarrow\mathcal B} = P_{\mathcal C}^{-1} P_{\mathcal B}.
\]

\revisit With $\mathcal B=\{(1,1),(1,2)\}$ and $\mathcal C=\{(1,0),(2,1)\}$:
\[
  \aug{cc|cc}{1&2&1&1\\0&1&1&2}
  \xrightarrow{R_1-2R_2}
  \aug{cc|cc}{1&0&-1&-3\\0&1&1&2},
  \qquad
  P_{\mathcal C\leftarrow\mathcal B} = \begin{bmatrix}-1&-3\\1&2\end{bmatrix}, \qquad
  P_{\mathcal B\leftarrow\mathcal C} = (P_{\mathcal C\leftarrow\mathcal B})^{-1} = \begin{bmatrix}2&3\\-1&-1\end{bmatrix}.
\]
Thus
\[
  [\vv x]_{\mathcal C} = \begin{bmatrix}-1&-3\\1&2\end{bmatrix}\begin{bmatrix}2\\-1\end{bmatrix} = \begin{bmatrix}1\\0\end{bmatrix},
\]
meaning $\vv x=2\vv b_1-\vv b_2=(1,0)=1\vv c_1+0\vv c_2$.

\begin{fig}
\begin{tikzpicture}[>=Stealth]
  % (a) Two coordinate representations of the same vector x = (1,0)
  \begin{scope}[xshift=0cm,yshift=0cm,scale=1.22]
    \draw[ink line,figaxis] (-0.5,0) -- (2.85,0) node[right,font=\scriptsize] {\InkLabel{$x_1$}};
    \draw[ink line,figaxis] (0,-0.45) -- (0,2.4) node[above,font=\scriptsize] {\InkLabel{$x_2$}};

    % B-basis (single restrained accent figblue) and dashed path 2b1 - b2
    \InkStroke[figvec,draw=figblue,line width=0.65pt]{(0,0) -- (1,1)}
    \node[font=\scriptsize,text=figblue!90!black] at (0.82,1.10) {\InkLabel{$\vv b_1$}};
    \draw[ink line,figvec,draw=figblue,line width=0.65pt] (0,0) -- (1,2) node[above left=1pt,font=\scriptsize,text=figblue!90!black] {\InkLabel{$\vv b_2$}};
    \draw[ink guide arrow,draw=figblue!75] (1,1) -- (2,2) node[above=2pt,font=\scriptsize,text=figblue!90!black] {\InkLabel{$2\vv b_1$}};
    \draw[ink guide arrow,draw=figblue!75,shorten >=2.8pt] (2,2) -- (1,0) node[pos=0.35,right=3pt,font=\scriptsize,text=figblue!90!black] {\InkLabel{$-\vv b_2$}};

    % C-basis vector c2 and focal vector x = c1 (pure black ink)
    \draw[ink line,figvec,draw=black,line width=0.65pt] (0,0) -- (2,1) node[right=3pt,font=\scriptsize,text=black!75] {\InkLabel{$\vv c_2$}};
    \InkStroke[figvec,draw=black,line width=0.65pt,shorten >=2.4pt]{(0,0) -- (1,0)}
    \node[figdot] at (1,0) {};
    \node[below=3pt,font=\scriptsize] at (1,0) {\InkLabel{$\vv x=\vv c_1$}};

    \node[font=\footnotesize] at (1.15,-0.78) {\InkLabel{(a)}};
  \end{scope}

  % (b) Plain commutative triangle on white (no tinted boxes)
  \begin{scope}[xshift=6.4cm,yshift=0.55cm]
    \node[font=\small] (std) at (2.25,1.85) {\InkLabel{$\vv x\in\RR^n$}};
    \node[font=\small] (B) at (0,0) {\InkLabel{$[\vv x]_{\mathcal B}$}};
    \node[font=\small] (C) at (4.5,0) {\InkLabel{$[\vv x]_{\mathcal C}$}};

    \draw[ink line,-{Stealth[length=4.0pt,width=2.3pt]},black!85,line width=0.55pt,shorten <=4pt,shorten >=4pt] (B) -- node[above left=1pt,font=\scriptsize] {\InkLabel{$P_{\mathcal B}$}} (std);
    \draw[ink line,-{Stealth[length=4.0pt,width=2.3pt]},black!85,line width=0.55pt,shorten <=4pt,shorten >=4pt] (C) -- node[above right=1pt,font=\scriptsize] {\InkLabel{$P_{\mathcal C}$}} (std);
    \draw[ink line,-{Stealth[length=4.2pt,width=2.4pt]},black,line width=0.65pt,shorten <=4pt,shorten >=4pt] (B) -- node[below=4pt,font=\scriptsize] {\InkLabel{$P_{\mathcal C\leftarrow\mathcal B}=P_{\mathcal C}^{-1}P_{\mathcal B}$}} (C);

    \node[font=\footnotesize] at (2.25,-1.50) {\InkLabel{(b)}};
  \end{scope}
\end{tikzpicture}
\figcap{\textbf{(a)} Both $\mathcal B$-coordinates $(2,-1)$ (blue) and $\mathcal C$-coordinates $(1,0)$ locate the same vector $\vv x=(1,0)=2\vv b_1-\vv b_2=\vv c_1$. \textbf{(b)} Converting $\mathcal B\to\mathcal C$ in $\RR^n$ factors through standard coordinates as $P_{\mathcal C}^{-1}P_{\mathcal B}$.}
\end{fig}

\begin{example*}[A change-of-coordinates matrix for abstract bases]
In $V$ with $\vv b_1=2\vv c_1-\vv c_2$ and $\vv b_2=\vv c_1+3\vv c_2$: $P_{\mathcal C\leftarrow\mathcal B}=\begin{psmallmatrix}2&1\\-1&3\end{psmallmatrix}$, so
    \[
      [3\vv b_1+\vv b_2]_{\mathcal C}=\begin{bmatrix}2&1\\-1&3\end{bmatrix}\begin{bmatrix}3\\1\end{bmatrix}=\begin{bmatrix}7\\0\end{bmatrix} \quad (=7\vv c_1).
    \]
Check directly: $3(2\vv c_1-\vv c_2)+(\vv c_1+3\vv c_2)=7\vv c_1$.
\end{example*}

\begin{example*}[A change-of-coordinates matrix in $\mathbb P_2$]
In $\mathbb P_2$ with $\mathcal B=\{1,1+t,1+t+t^2\}$ and $\mathcal C=\{1,t,t^2\}$:
    \[
      P_{\mathcal C\leftarrow\mathcal B} = \begin{bmatrix}1&1&1\\0&1&1\\0&0&1\end{bmatrix},\qquad
      P_{\mathcal B\leftarrow\mathcal C} = \begin{bmatrix}1&-1&0\\0&1&-1\\0&0&1\end{bmatrix},
    \]
so
\[
  [2+3t+4t^2]_{\mathcal B} = P_{\mathcal B\leftarrow\mathcal C}\begin{bmatrix}2\\3\\4\end{bmatrix} = \begin{bmatrix}-1\\-1\\4\end{bmatrix}.
\] The columns of $P_{\mathcal C\leftarrow\mathcal B}$ are the coefficient lists of $1$, $1+t$, $1+t+t^2$. Check: $-1\cdot1-1\cdot(1+t)+4(1+t+t^2)=2+3t+4t^2$.
\end{example*}

\begin{shortcuts}
  \ii \textbf{Change of basis in $\RR^n$:} Put the \emph{target} basis $\mathcal C$ on the left: $[\,\mathcal C\mid\mathcal B\,]\to [\,I\mid P_{\mathcal C\leftarrow\mathcal B}\,]$.
  \ii \textbf{Reverse direction:} Invert $P_{\mathcal C\leftarrow\mathcal B}$ rather than reducing a second system.
\end{shortcuts}

\begin{trapbox}
  \ii \textbf{Left--right order:} Reducing $[\,\mathcal B\mid\mathcal C\,]$ produces $P_{\mathcal B\leftarrow\mathcal C}$ (the inverse of $P_{\mathcal C\leftarrow\mathcal B}$).
  \ii \textbf{Columns vs.\ rows:} The coordinates of $\vv b_1$ in $\mathcal C$ form \emph{column~$1$} of $P_{\mathcal C\leftarrow\mathcal B}$, never row~$1$.
\end{trapbox}

\subsection{Eigenvectors and eigenvalues (Lay 5.1)}

\begin{qbox}
  The matrix $A=\begin{psmallmatrix}1&6\\5&2\end{psmallmatrix}$ sends the vector $(3,-2)$ off its own line to $(-9,11)$, yet sends $(6,-5)$ to $(-24,20)=-4(6,-5)$ along the \emph{same} line. Which vectors does a matrix $A$ merely stretch or flip, and how do they compute $A^k\vv x$ for large $k$ without multiplying matrices $k$ times?
\end{qbox}

\noindent\textbf{Definition (Lay 5.1).} An \vocab{eigenvector} of an $n\times n$ matrix $A$ is a \emph{nonzero} vector $\vv x\in\RR^n$ such that $A\vv x=\lambda\vv x$ for some scalar $\lambda$. A scalar $\lambda$ is an \vocab{eigenvalue} of $A$ if $A\vv x=\lambda\vv x$ has a nontrivial solution; any such $\vv x\neq\vv 0$ is an eigenvector corresponding to $\lambda$. The set of all solutions of $(A-\lambda I)\vv x=\vv 0$, namely $\Nul(A-\lambda I)$, is the \vocab{eigenspace} of $A$ for $\lambda$.

Eigenvectors make powers of $A$ easy to compute: if $A\vv v=\lambda\vv v$, then $A^2\vv v=A(\lambda\vv v)=\lambda^2\vv v$, and by induction $A^k\vv v=\lambda^k\vv v$. Expanding any starting vector $\vv x_0 = c_1\vv v_1+c_2\vv v_2$ solves the discrete dynamical system $\vv x_{k+1}=A\vv x_k$ directly:
\[
  \vv x_k = A^k\vv x_0 = c_1\lambda_1^k\vv v_1 + c_2\lambda_2^k\vv v_2.
\]

\method{How to check eigenvectors and eigenvalues, and find an eigenspace.}
\begin{steps}
  \ii \emph{Is $\vv u\ne\vv 0$ an eigenvector?} Compute $A\vv u$ and check whether it is a multiple $\lambda\vv u$; the multiple is the eigenvalue.
  \ii \emph{Is $\lambda$ an eigenvalue?} Row reduce $A-\lambda I$; $\lambda$ is an eigenvalue exactly when there is a free variable.
  \ii \emph{Eigenspace:} write the solutions of $(A-\lambda I)\vv x=\vv 0$ in parametric vector form; the vectors attached to the free variables form a basis of $\Nul(A-\lambda I)$.
\end{steps}
\emph{Edge cases.} $\vv 0$ is never an eigenvector, but $\lambda=0$ can be an eigenvalue (exactly when $A$ is not invertible). To make an eigenspace have a given dimension, choose a parameter so that $A-\lambda I$ has that many free variables.

\begin{example*}[A number that is not an eigenvalue]
For $A=\begin{psmallmatrix}4&2\\1&3\end{psmallmatrix}$ and $\lambda=3$, $A-3I=\begin{psmallmatrix}1&2\\1&0\end{psmallmatrix}$ has two pivots and no free variable, so $3$ is not an eigenvalue. The vector $(2,1)$ is an eigenvector: $A(2,1)=(10,5)=5(2,1)$, with eigenvalue $5$.
\end{example*}

\begin{example*}[An eigenspace that is a plane]
  For $A$ below, is $\lambda=2$ an eigenvalue and what is $E(2,A)=\Nul(A-2I)$?
  \[
    A=\begin{bmatrix}4&-1&6\\2&1&6\\2&-1&8\end{bmatrix},\qquad
    A-2I = \begin{bmatrix}2&-1&6\\2&-1&6\\2&-1&6\end{bmatrix}
    \xrightarrow{R_2-R_1,\ R_3-R_1}
    \begin{bmatrix}2&-1&6\\0&0&0\\0&0&0\end{bmatrix}.
  \]
  Since there are two free variables ($x_2,x_3$), $\lambda=2$ \emph{is} an eigenvalue, and $x_1=\frac12 x_2-3x_3$ gives eigenspace basis $\{(1,2,0),(-3,0,1)\}$. On this entire plane, $A$ acts purely as multiplication by $2$.
\end{example*}

\begin{example*}[Choosing a parameter to get a two-dimensional eigenspace]
For the matrix $A$ below,
\[
  A=\begin{bmatrix}5&h&1\\0&5&0\\0&0&3\end{bmatrix},\qquad A-5I=\begin{bmatrix}0&h&1\\0&0&0\\0&0&-2\end{bmatrix}.
\]
If $h\ne0$, rows $1$ and $3$ give pivots in columns $2$ and $3$: one free variable, a line. If $h=0$, both nonzero rows have their pivot in column $3$, so there is one pivot and two free variables. The eigenspace for $5$ is two-dimensional exactly when $h=0$.
\end{example*}

\begin{theorem*}[Eigenvalues of a triangular matrix (Lay 5.1, Thm.~1)]\target{thm-lay-5-1}
  The eigenvalues of a triangular matrix $A$ are the entries on its main diagonal.
\end{theorem*}
\begin{proof}
  If $A$ is upper- or lower-triangular, so is $A-\lambda I$, with diagonal entries $a_{11}-\lambda,\dots,a_{nn}-\lambda$. Now $\lambda$ is an eigenvalue $\iff (A-\lambda I)\vv x=\vv 0$ has a nontrivial solution $\iff A-\lambda I$ is not invertible (\ilink{fact-w2}{W2}) $\iff \det(A-\lambda I)=\prod_{i=1}^n (a_{ii}-\lambda)=0$ (\ilink{fact-w3}{W3}) $\iff \lambda=a_{ii}$ for some $i$.
\end{proof}

Notice that $\lambda=0$ is an eigenvalue of $A \iff A\vv x=\vv 0$ has a nontrivial solution $\iff A$ is \emph{not} invertible (this updates the \vocab{Invertible Matrix Theorem}: $A$ is invertible $\iff 0$ is not an eigenvalue of $A$).

\begin{theorem*}[Linear independence of eigenvectors (Lay 5.1, Thm.~2)]\target{thm-lay-5-2}
  If $\vv v_1,\dots,\vv v_r$ are eigenvectors corresponding to \textbf{distinct} eigenvalues $\lambda_1,\dots,\lambda_r$ of an $n\times n$ matrix $A$, then $\{\vv v_1,\dots,\vv v_r\}$ is linearly independent.
\end{theorem*}
\begin{ideabox}
  Think of two vectors first: if $\vv v_2=c\vv v_1$, applying $A$ gives $\lambda_2\vv v_2=c\lambda_1\vv v_1$, whereas multiplying by $\lambda_2$ gives $\lambda_2\vv v_2=c\lambda_2\vv v_1$. Subtracting leaves $c(\lambda_1-\lambda_2)\vv v_1=\vv 0$, forcing $c=0$ since $\lambda_1\ne\lambda_2$. In general, apply the operator $A-\lambda_{p+1}I$ to the \emph{first} vector $\vv v_{p+1}$ that depends on earlier ones: it annihilates $\vv v_{p+1}$ and multiplies each remaining term by $\lambda_i-\lambda_{p+1}\neq0$.
\end{ideabox}
\begin{proof}
  \begin{steps}
    \ii Suppose for contradiction that $\{\vv v_1,\dots,\vv v_r\}$ is linearly dependent. Since $\vv v_1\ne\vv 0$, by (\ilink{fact-w6}{W6}) there is a first vector $\vv v_{p+1}$ ($1\le p<r$) that is a linear combination of the preceding vectors $\{\vv v_1,\dots,\vv v_p\}$ (which are linearly independent):
    \[
      \vv v_{p+1} = c_1\vv v_1 + \dots + c_p\vv v_p.
    \]
    \ii Apply $A$ to both sides (using $A\vv v_i=\lambda_i\vv v_i$): $\lambda_{p+1}\vv v_{p+1}=c_1\lambda_1\vv v_1+\dots+c_p\lambda_p\vv v_p$.
    \ii Multiply step~1 by $\lambda_{p+1}$: $\lambda_{p+1}\vv v_{p+1}=c_1\lambda_{p+1}\vv v_1+\dots+c_p\lambda_{p+1}\vv v_p$.
    \ii Subtracting eliminates $\vv v_{p+1}$: $\vv 0 = c_1(\lambda_1-\lambda_{p+1})\vv v_1+\dots+c_p(\lambda_p-\lambda_{p+1})\vv v_p$. By independence of $\vv v_1,\dots,\vv v_p$ and $\lambda_i\ne\lambda_{p+1}$, every $c_i=0$, forcing $\vv v_{p+1}=\vv 0$, impossible for an eigenvector.
  \end{steps}
\end{proof}

As an immediate consequence, an $n\times n$ matrix has at most $n$ distinct eigenvalues (otherwise it would have $n+1$ independent vectors in $\RR^n$, impossible by \ilink{fact-w1}{W1}).

\begin{fig}
\begin{tikzpicture}[>=Stealth]
  % (a) Eigenvectors stay on their invariant lines (inputs black, images figblue)
  \begin{scope}[xshift=0cm,yshift=0cm,scale=0.145]
    \draw[ink line,figaxis] (-26,0) -- (15,0) node[right,font=\scriptsize] {\InkLabel{$x_1$}};
    \draw[ink line,figaxis] (0,-9.5) -- (0,23) node[above,font=\scriptsize] {\InkLabel{$x_2$}};

    % Dotted invariant lines
    \draw[ink guide,draw=black!45] (-25.8,21.5) -- (9.6,-8.0);
    \draw[ink guide,draw=black!45] (-5.5,-5.5) -- (14.2,14.2);

    % Images under A (restrained figblue) drawn first so inputs sit crisply on top
    \draw[ink line,figvec,draw=figblue,line width=0.65pt] (0,0) -- (-24,20) node[above right=2pt,font=\scriptsize,text=figblue!90!black] {\InkLabel{$A\vv v=-4\vv v$}};
    \draw[ink line,figvec,draw=figblue,line width=0.65pt] (0,0) -- (12.6,12.6) node[above left=2pt,font=\scriptsize,text=figblue!90!black] {\InkLabel{$A\vv w=7\vv w$}};

    % Input eigenvectors (pure black)
    \draw[ink line,figvec,draw=black,line width=0.65pt] (0,0) -- (6,-5) node[below=3pt,font=\scriptsize] {\InkLabel{$\vv v$}};
    \InkStroke[figvec,draw=black,line width=0.65pt]{(0,0) -- (1.8,1.8)}
    \node[right=4pt,font=\scriptsize] at (2.1,2.0) {\InkLabel{$\vv w$}};

    \node[font=\footnotesize] at (-5.5,-13.0) {\InkLabel{(a)}};
  \end{scope}

  % (b) Non-eigenvector leaves its line
  \begin{scope}[xshift=6.6cm,yshift=-0.25cm,scale=0.195]
    \draw[ink line,figaxis] (-11.5,0) -- (7.0,0) node[right,font=\scriptsize] {\InkLabel{$x_1$}};
    \draw[ink line,figaxis] (0,-5.5) -- (0,13.5) node[above,font=\scriptsize] {\InkLabel{$x_2$}};

    % Dotted line spanned by u = (3,-2)
    \draw[ink guide,draw=black!45] (-10.8,7.2) -- (4.8,-3.2);

    % Input vector u (black) and image Au (figblue, solid)
    \draw[ink line,figvec,draw=black,line width=0.65pt] (0,0) -- (3,-2) node[below=3pt,font=\scriptsize] {\InkLabel{$\vv u$}};
    \draw[ink line,figvec,draw=figblue,line width=0.65pt] (0,0) -- (-9,11) node[above left=2pt,font=\scriptsize,text=figblue!90!black] {\InkLabel{$A\vv u$}};

    \node[font=\footnotesize] at (-2.0,-8.4) {\InkLabel{(b)}};
  \end{scope}
\end{tikzpicture}
\figcap{Action of $A=\begin{psmallmatrix}1&6\\5&2\end{psmallmatrix}$ (inputs black, images blue): \textbf{(a)} eigenvectors $\vv v,\vv w$ stay on their dotted lines; \textbf{(b)} $\vv u=(3,-2)$ is sent off its line to $A\vv u=(-9,11)$.}
\end{fig}

\revisit The invariant directions of $A$ are its eigenvectors: $A(1,1)=(7,7)=7(1,1)$ and $A(6,-5)=-4(6,-5)$. Writing $(7,-4)^{\mathsf T}=(1,1)^{\mathsf T}+(6,-5)^{\mathsf T}$ determines every power directly:
\[
  A^k\begin{bmatrix}7\\-4\end{bmatrix} = 7^k\begin{bmatrix}1\\1\end{bmatrix} + (-4)^k\begin{bmatrix}6\\-5\end{bmatrix}.
\]

\begin{shortcuts}
  \ii \textbf{Checking an eigenvector $\vv u$ vs.\ eigenvalue $\lambda$:} For $\vv u$, compute $A\vv u$ and check if it is $\lambda\vv u$; for $\lambda$, row reduce $A-\lambda I$ and check for a free variable.
  \ii \textbf{Equal row sums:} If every row of $A$ sums to $s$, then $s$ is an eigenvalue with eigenvector $(1,1,\dots,1)^{\mathsf T}$ (in the question, rows add to $7$).
  \ii \textbf{Powers \& inverses:} If $A\vv x=\lambda\vv x$ ($\vv x\neq\vv 0$), then $A^k\vv x=\lambda^k\vv x$, and if $A$ is invertible, $\lambda\ne 0$ and $A^{-1}\vv x=\lambda^{-1}\vv x$.
\end{shortcuts}

\begin{trapbox}
  \ii \textbf{The zero vector is never an eigenvector:} $A\vv 0=\lambda\vv 0$ holds trivially for every $\lambda$, whereas $0$ \emph{can} be an eigenvalue.
  \ii \textbf{Row-reducing $A$ changes eigenvalues:} $A=\begin{psmallmatrix}1&6\\5&2\end{psmallmatrix}$ has eigenvalues $7,-4$, whereas $R_2-5R_1$ gives $\begin{psmallmatrix}1&6\\0&-28\end{psmallmatrix}$ with eigenvalues $1,-28$. Always row reduce $A-\lambda I$, never $A$.
  \ii \textbf{\ilink{thm-lay-5-2}{Thm.~5.2} goes one way:} Two independent eigenvectors can share the same eigenvalue.
\end{trapbox}

\subsection{The characteristic equation (Lay 5.2)}

\begin{qbox}
  Testing candidate values of $\lambda$ one by one is slow. How do we find \emph{all} eigenvalues of $A=\begin{psmallmatrix}3&1\\1&3\end{psmallmatrix}$ simultaneously as roots of a single polynomial equation, and why do similar matrices always share the same eigenvalues?
\end{qbox}

\begin{theorem*}[Properties of determinants recap (Lay 5.2, Thm.~3) \& Characteristic Equation]\target{thm-lay-5-3}\target{name-characteristic-equation}
  Recall for $n\times n$ matrices $A,B$ (\ilink{fact-w2}{W2}, \ilink{fact-w3}{W3}): $A$ is invertible $\iff \det A\ne 0$; $\det(AB)=\det A\det B$; $\det(A^{\mathsf T})=\det A$; triangular determinants equal the product of diagonal entries; and row operations scale/flip/preserve $\det$. Consequently, a scalar $\lambda$ is an eigenvalue of $A$ if and only if $\lambda$ satisfies the \textbf{characteristic equation}
  \[
    \det(A - \lambda I) = 0.
  \]
\end{theorem*}
\begin{proof}
  $A\vv x=\lambda\vv x$ has a nontrivial solution $\iff (A-\lambda I)\vv x=\vv 0$ has a nontrivial solution $\iff A-\lambda I$ is not invertible (\ilink{fact-w2}{W2}) $\iff \det(A-\lambda I)=0$ (\ilink{fact-w3}{W3}).
\end{proof}

\begin{proposition*}[The characteristic polynomial has degree $n$]\target{prop-charpoly-degree}
  For an $n\times n$ matrix $A$, $\det(A-\lambda I)$ is a polynomial in $\lambda$ of degree $n$:
  \[
    \det(A-\lambda I)=(-1)^n\lambda^n+(-1)^{n-1}(a_{11}+\dots+a_{nn})\lambda^{n-1}+(\text{terms of degree}\le n-2).
  \]
\end{proposition*}
\begin{ideabox}
  Only the diagonal entries of $A-\lambda I$ contain $\lambda$. Expanding along a row multiplies one entry by a smaller determinant, so each $\lambda$ can raise the degree by at most one. The product of the diagonal entries uses all $n$ of them; every other term of the expansion skips at least two diagonal entries.
\end{ideabox}
\begin{proof}
  \begin{steps}
    \ii \textbf{Counting lemma.} If a square matrix has numbers as entries except for $k$ entries of the form $c-\lambda$, no two in the same row, then its determinant is a polynomial in $\lambda$ of degree at most $k$. By induction on the size: expand along a row (\ilink{fact-w3}{W3}). If the row has no $\lambda$-entry, each term is a number times a minor with at most $k$ $\lambda$-entries. If the row has the $\lambda$-entry, that term is $(c-\lambda)$ times a minor with at most $k-1$ $\lambda$-entries (degree $\le 1+(k-1)$), and the other terms are numbers times minors with at most $k-1$. In all cases the degree is at most $k$.
    \ii \textbf{Expansion along row~1.} $\det(A-\lambda I)=(a_{11}-\lambda)\det M_{11}+\sum_{j\ge2}(\pm a_{1j})\det M_{1j}$, where $M_{1j}$ deletes row $1$ and column $j$ of $A-\lambda I$. For $j\ge2$, $M_{1j}$ has lost the diagonal entries in row $1$ and in column $j$, so it has at most $n-2$ $\lambda$-entries and $\det M_{1j}$ has degree at most $n-2$ by Step~1.
    \ii \textbf{Induction.} $M_{11}=A_{11}-\lambda I_{n-1}$. If the formula holds for $(n-1)\times(n-1)$ matrices, $\det M_{11}=(-1)^{n-1}\lambda^{n-1}+(-1)^{n-2}(a_{22}+\dots+a_{nn})\lambda^{n-2}+\dots$. Multiplying by $a_{11}-\lambda$ gives $(-1)^n\lambda^n+(-1)^{n-1}(a_{11}+a_{22}+\dots+a_{nn})\lambda^{n-1}+(\text{degree}\le n-2)$. Adding the terms of Step~2 does not change the two leading terms. The case $n=1$ is $a_{11}-\lambda$.
  \end{steps}
\end{proof}

The polynomial $\det(A-\lambda I)$ is the \vocab{characteristic polynomial} of $A$, and its roots are the eigenvalues. The \vocab{(algebraic) multiplicity} of an eigenvalue $\lambda$ is its multiplicity as a root of $\det(A-\lambda I)$. Setting $\lambda=0$ gives $\det(A-0I)=\det A$. If the characteristic polynomial factors completely into linear factors, $\det(A-\lambda I)=(\lambda_1-\lambda)\cdots(\lambda_n-\lambda)$ (the leading coefficient is $(-1)^n$ by the proposition), so $\det A=\lambda_1\cdots\lambda_n$, the product of the eigenvalues counted with multiplicity, and comparing the $\lambda^{n-1}$ coefficients gives $\lambda_1+\dots+\lambda_n=a_{11}+\dots+a_{nn}=\tr A$. Moreover, since $\det(A^{\mathsf T}-\lambda I)=\det\bigl((A-\lambda I)^{\mathsf T}\bigr)=\det(A-\lambda I)$ by (\ilink{fact-w3}{W3}), $A$ and $A^{\mathsf T}$ always have the \emph{same eigenvalues} (though generally different eigenvectors).

\revisit[, with the $2\times 2$ formula] For any $2\times 2$ matrix $A=\begin{psmallmatrix}a&b\\c&d\end{psmallmatrix}$ with $\tr A=a+d$:
\[
  \det(A-\lambda I) = (a-\lambda)(d-\lambda)-bc = \lambda^2-(\tr A)\lambda+\det A.
\]
If this has roots $\lambda_1,\lambda_2$, comparing with $(\lambda-\lambda_1)(\lambda-\lambda_2)$ gives $\lambda_1+\lambda_2=\tr A$ and $\lambda_1\lambda_2=\det A$. For $A=\begin{psmallmatrix}3&1\\1&3\end{psmallmatrix}$, $\tr A=6$, $\det A=8$, so $\lambda^2-6\lambda+8=(\lambda-2)(\lambda-4)=0$ gives eigenvalues $2,4$ with eigenvectors $(-1,1)$ and $(1,1)$.

\begin{example*}[A $3\times 3$ characteristic polynomial by expanding down a column]
  Expand down column~$1$, which has two zeros, and keep linear factors intact instead of multiplying out:
  \[
    A=\begin{bmatrix}1&2&1\\0&3&0\\0&1&2\end{bmatrix}, \qquad
    \det(A-\lambda I) = (1-\lambda)\det\begin{bmatrix}3-\lambda&0\\1&2-\lambda\end{bmatrix} = (1-\lambda)(3-\lambda)(2-\lambda),
  \]
  giving eigenvalues $1,2,3$ immediately.
\end{example*}

\method{How to find all eigenvalues and their multiplicities.}
\begin{steps}
  \ii For a triangular matrix, read the diagonal (\ilink{thm-lay-5-1}{Thm.~5.1}).
  \ii For a $2\times2$ matrix, solve $\lambda^2-(\tr A)\lambda+\det A=0$.
  \ii Otherwise form $A-\lambda I$ and expand $\det(A-\lambda I)$ along the row or column with the most zeros, keeping factors like $(2-\lambda)$ unmultiplied.
  \ii Factor; each root is an eigenvalue, and the power of its factor is its multiplicity. Check: the eigenvalues add up to $\tr A$ and multiply to $\det A$ (when the polynomial factors completely).
\end{steps}
\emph{Edge cases.} A repeated root has multiplicity at least $2$. A quadratic factor with negative discriminant has no real roots; those eigenvalues are complex. The eigenvalues of $A$ are not those of its echelon form.

\begin{example*}[A repeated eigenvalue]
For the matrix $A$ below, row~1 of $A-\lambda I$ is $(2-\lambda,0,0)$, so
\[
  A=\begin{bmatrix}2&0&0\\1&2&1\\-1&0&3\end{bmatrix},\qquad
  \det(A-\lambda I)=(2-\lambda)\det\begin{bmatrix}2-\lambda&1\\0&3-\lambda\end{bmatrix}=(2-\lambda)^2(3-\lambda).
\]
The eigenvalue $2$ has multiplicity $2$ and $3$ has multiplicity $1$. Check: $2+2+3=7=\tr A$ and $2\cdot2\cdot3=12=\det A$.
\end{example*}

\noindent\textbf{Definition.} Two $n\times n$ matrices $A$ and $B$ are \vocab{similar} if $B=P^{-1}AP$ for some invertible matrix $P$.

\begin{theorem*}[Similarity and characteristic polynomials (Lay 5.2, Thm.~4)]\target{thm-lay-5-4}
  If $n\times n$ matrices $A$ and $B$ are similar ($B=P^{-1}AP$), then they have the same characteristic polynomial, hence the exact same eigenvalues with the exact same multiplicities.
\end{theorem*}
\begin{proof}
  Notice $B-\lambda I = P^{-1}AP - \lambda P^{-1}P = P^{-1}(A-\lambda I)P$. Using (\ilink{fact-w3}{W3}),
  \[
    \det(B-\lambda I) = \det(P^{-1})\det(A-\lambda I)\det P = \det(P^{-1}P)\det(A-\lambda I) = \det(A-\lambda I).
  \]
\end{proof}

\revisit[ (similar matrices)] If $B=P^{-1}AP$, then $B-\lambda I=P^{-1}(A-\lambda I)P$, and the determinant of this product is $\det(A-\lambda I)$ because $\det P^{-1}\det P=1$. So $A$ and $B$ have the same characteristic polynomial and therefore the same eigenvalues with the same multiplicities (\ilink{thm-lay-5-4}{Thm.~5.4}). For example, $P=\begin{psmallmatrix}1&1\\0&1\end{psmallmatrix}$ and $A=\begin{psmallmatrix}3&1\\1&3\end{psmallmatrix}$ give $AP=\begin{psmallmatrix}3&4\\1&4\end{psmallmatrix}$ and $B=P^{-1}AP=\begin{psmallmatrix}2&0\\1&4\end{psmallmatrix}$, with $\tr B=6$ and $\det B=8$: eigenvalues $2,4$, as for $A$.

\begin{shortcuts}
  \ii \textbf{$2\times 2$ eigenvalues:} Solve $\lambda^2-(\tr A)\lambda+\det A=0$; verify $\lambda_1+\lambda_2=\tr A$ and $\lambda_1\lambda_2=\det A$.
  \ii \textbf{$2\times 2$ eigenvector shortcut:} If row~$1$ of $A-\lambda I$ is $(a-\lambda,b)\ne(0,0)$, an eigenvector is $(b,\lambda-a)$ (e.g.\ $b=0$ gives $(1,0)$, or inspect row~$2$).
  \ii \textbf{$3\times 3$ characteristic polynomial:} Expand along the row or column with the most zeros and factor rather than expanding a cubic.
\end{shortcuts}

\begin{trapbox}
  \ii \textbf{$\det(A-\lambda I)=0$, \emph{not} $\det A-\lambda=0$:} For $\begin{psmallmatrix}3&1\\1&3\end{psmallmatrix}$, $\det A-\lambda=8-\lambda$ wrongly predicts $8$.
  \ii \textbf{Same characteristic polynomial $\not\Rightarrow$ similar:} $\begin{psmallmatrix}2&1\\0&2\end{psmallmatrix}$ and $2I$ both have $(\lambda-2)^2$, yet $P^{-1}(2I)P=2I$ for every $P$.
  \ii \textbf{Similarity $\neq$ row equivalence:} Both matrices above row-reduce to $I$, yet are not similar.
\end{trapbox}

\section{Homework this week}
HW6, Sections 4.4 and 4.6, and HW7, Sections 5.1--5.2.
\begin{itemize}
  \ii \textbf{Section 4.4:} \#1, \#3 (use $\vv x=P_{\mathcal B}[\vv x]_{\mathcal B}$); \#7 (row reduce $[\,\vv b_1\ \vv b_2\ \vv b_3\mid \vv x\,]$); \#13 (equate coefficients of $1,t,t^2$, see P2(a)); \#15, \#17, \#19 (definition of $[\vv x]_{\mathcal B}=P_{\mathcal B}^{-1}\vv x$; $\dim\mathbb P_3=4$); \#25 ($P_{\mathcal B}^{-1}$); \#36 (convert to $\RR^n$ coordinates, apply Basis Thm., then solve weights).
  \ii \textbf{Section 4.6:} \#1, \#5, \#6 (abstract bases example); \#9 (row reduce $[\,\vv c_1\ \vv c_2\mid \vv b_1\ \vv b_2\,]$, then invert); \#12, \#13 (with $\mathcal C=\mathcal E$, $P_{\mathcal E\leftarrow\mathcal B}=P_{\mathcal B}$); \#15 ($\mathbb P_2$ coordinate columns).
  \ii \textbf{Section 5.1:} \#3, \#5 (multiply $A\vv u$ and compare with $\vv u$); \#7 (row reduce $A-4I$ for free variables); \#9, \#13 (parametric vector form of $(A-\lambda I)\vv x=\vv 0$); \#23, \#24, \#27, \#28, \#29, \#30 ($0$ eigenvalue $\iff$ singular; $A\vv x=\lambda\vv x$; \ilink{thm-lay-5-1}{Thm.~5.1} needs triangular; $A(2\vv v)=2(A\vv v)$); \#34 (apply $A$ twice to $A\vv x=\lambda\vv x$).
  \ii \textbf{Section 5.2:} \#5, \#7 ($\lambda^2-(\tr A)\lambda+\det A$); \#11, \#13 (expand along sparse row/column); \#15 (triangular); \#18 (choose $h$ so $A-5I$ has two free variables); \#19 (set $\lambda=0$); \#24, \#25 (\ilink{thm-lay-5-4}{Thm.~5.4} and the characteristic equation).
\end{itemize}

\section{Problem set}

\begin{prob}\target{prob-P1}
  \textbf{P1.}\ Find $[\vv x]_{\mathcal B}$ for $\vv x=(4,9)$ and $\mathcal B=\{(1,3),(2,5)\}$.
\end{prob}

\begin{prob}\target{prob-P2}
  \textbf{P2.}\ \begin{enumerate}[(a)]
    \ii $\mathcal B=\{1+t,\ t+t^2,\ 1+t^2\}$ is a basis of $\mathbb P_2$. Find $[p]_{\mathcal B}$ for $p=3+2t-t^2$.
    \ii Use coordinate vectors to decide whether $1-t,\ t-t^2,\ 1-t^2$ are linearly independent.
  \end{enumerate}
\end{prob}

\begin{prob}\target{prob-P3}
  \textbf{P3.}\ $\mathcal B=\{(1,2),(3,4)\}$ and $\mathcal C=\{(1,1),(0,1)\}$. Find $P_{\mathcal C\leftarrow\mathcal B}$ and $P_{\mathcal B\leftarrow\mathcal C}$. If $[\vv x]_{\mathcal B}=(2,-1)$, find $[\vv x]_{\mathcal C}$.
\end{prob}

\begin{prob}\target{prob-P4}
  \textbf{P4.}\ For the matrix $A$ below, is $\vv v=(1,-1,0)$ an eigenvector? If so, for which eigenvalue? Is $4$ an eigenvalue?
  \[
    A = \begin{bmatrix}3&0&1\\-1&2&-1\\0&0&1\end{bmatrix}.
  \]
\end{prob}

\begin{prob}\target{prob-P5}
  \textbf{P5.}\ For the same $A$, find the characteristic polynomial, all eigenvalues, and a basis of the eigenspace for $\lambda=1$.
\end{prob}

\begin{prob}\target{prob-P6}
  \textbf{P6.}\ Find the characteristic polynomials and eigenvalues of the matrices below:
  \[
    \begin{bmatrix}2&7\\7&2\end{bmatrix},\qquad \begin{bmatrix}5&-3\\-4&3\end{bmatrix}.
  \]
\end{prob}

\begin{prob}\target{prob-P7}
  \textbf{P7.}\ List the eigenvalues, with multiplicities, of the matrix below. Is it invertible?
  \[
    \begin{bmatrix}2&3&-1&4\\0&-1&5&2\\0&0&2&7\\0&0&0&0\end{bmatrix}.
  \]
\end{prob}

\begin{prob}\target{prob-P8}
  \textbf{P8.}\ Suppose $A\vv x=\lambda\vv x$ with $\vv x\ne\vv 0$.
  \begin{enumerate}[(a)]
    \ii Show that $\vv x$ is an eigenvector of $A^2$ and of $A+3I$; for which eigenvalues?
    \ii If $A$ is invertible, show $\lambda\ne 0$ and $A^{-1}\vv x=\lambda^{-1}\vv x$.
  \end{enumerate}
\end{prob}

\end{document}
